\section{Square products with short selections}
\label{p13:short:section}

The conditional count in the next section must handle a relation
even when none of its selected integers has a large prime kernel.
For selections of at most a small multiple of \(\log\log N\),
we instead count square products through moments. The constants
in this argument are explicit in the selection size; a
fixed-degree result alone would not give the required range.

For a fixed number of unrestricted coordinates, global counts of
tuples whose product is a square were obtained by de~la~Bret\`eche,
Kurlberg, and Shparlinski
\cite[Theorem~2.2 and Corollary~2.4]{bretkurlbergshparlinski2021}.
Here the selections lie in short additive windows, and the ensuing
conditional estimates are uniform in arbitrary fixed source windows
even when the number of selected factors grows with $N$.

The quadratic kernel argument is related to the decomposition
in de la Bret\`eche, Wang, and Xu
\cite[Lemmas~6.2--6.3]{bretwngxu2026}. We give the estimates
needed here, including their dependence on the number and
location of the roots.


Write
\[
 P_A(X)=\prod_{a\in A}(X+a),\qquad
 A\subset[-DB,DB]\cap\Z,
 \qquad 2\le d=|A|\le c\log\log N,
\]
where $d$ is even and $D,c>0$ are fixed.  The interval of summation
has height and length comparable to $N$.

\begin{proposition}[Uniform moments for short split products]
\label{p13:short:moments}
For every fixed even integer \(v\ge2\), with the preceding notation,
\begin{equation}
 \E\left|\sum_{n\asymp N}f(P_A(n))\right|^v
 \le N^{v/2+c[v(v-1)+2]\log2+o(1)}.
 \label{p13:short:moment-bound}
\end{equation}
In particular, for each fixed \(v\) and \(\varepsilon>0\),
one may choose \(c>0\) so that the \(L^v\) norm is
\(O(N^{1/2+\varepsilon})\), uniformly in all the indicated
motifs. The summation interval may have any fixed comparable
length and positive height.
\end{proposition}

The proof, in Section~\ref{p14:comp:short:moments}, pairs the linear
factors into quadratics and controls the common kernel by their
resultant.  In the moment expansion, primes are assigned to the even
subsets of indices on which they occur.  Averaging the exposure order
then gives the exponent $v/2$.  The explicit term
$c[v(v-1)+2]\log2$ is what permits a growing selection size.

\begin{corollary}[Short projections in an old space]
\label{p13:short:old-space}
Fix \(r,s,D\) and \(\varepsilon>0\). Let \(V_{\rm old}\)
be generated by the square classes of all values in \(s\)
fixed blocks of diameter and cardinality \(O_D(B)\).
For a sufficiently small fixed \(c>0\), the number of
lists of \(r\) new centres admitting nonempty even
selections \(A_i\), \(|A_i|\le c\log\log N\), whose
square classes sum to an element of \(V_{\rm old}\), is
\(O_{r,s,D,C,\varepsilon}(N^{r/2+\varepsilon})\),
uniformly in the old blocks.
\end{corollary}

\begin{proof}
For a square class \(u\), let \(\chi_u(f)\) denote its
character, and put \(\Phi_V=\sum_{u\in V}\chi_u\), summing
over distinct classes. Character orthogonality gives
\[
 \Phi_V\ge0,\qquad \E\Phi_V=1,\qquad
 \|\Phi_V\|_a=2^{\dim V(1-1/a)}\quad(a\ge1).
\]
For fixed motifs, with \(X_i=\sum_z f(P_{A_i}(z))\),
the expectation \(\E(\Phi_{V_{\rm old}}\prod_iX_i)\)
counts exactly the required lists. For each list, a unique
old character cancels the new sum if it belongs to the old
space. All terms after expectation are zero or one.

Since \(\dim V_{\rm old}=O_{s,D}(B)\), choose a fixed
\(a>1\) sufficiently close to one that
\(\|\Phi_{V_{\rm old}}\|_a\le N^{\varepsilon'}\).
Next choose a fixed even \(v\ge ra/(a-1)\), and then choose
\(c\) by Proposition~\ref{p13:short:moments}.
H\"older gives
\[
 \left|\E\left(\Phi_{V_{\rm old}}\prod_iX_i\right)\right|
 \le N^{\varepsilon'}\prod_i
       \|X_i\|_{ra/(a-1)}
 \ll N^{r/2+(r+1)\varepsilon'}.
\]
The motifs have
\(\exp(O_{r,D}(c(\log\log N)^2))=N^{o(1)}\) choices.
Choose the internal margins in the stated \(\varepsilon\).
There is no enumeration of the \(2^{\dim V_{\rm old}}\)
old classes.
\end{proof}

\subsection{Short split products: kernels and moment enumeration}
\label{p14:comp:short}

This section proves the short-support moment input from
Section~\ref{p13:short:section}.  The degree is allowed to grow as a
sufficiently small multiple of $\log\log N$; all estimates below
retain that dependence.

\begin{lemma}[A quadratic kernel with a prescribed divisor]
\label{p13:short:quadratic}
Let \(1\le h\le2DB\), where \(D\) is fixed, and \(N\le n\le2N\).
Suppose \(N^\beta\le q=\sf(n(n+h))<2N^\beta\) and a fixed squarefree
integer \(e\) divides \(q\). Put \(\delta=\log e/\log N\).
Uniformly in these data, the number of \(n\) is at most
\[
 N^{c_2(\beta,\delta)+o(1)},\qquad
 c_2(\beta,\delta)=
 \min\{1,\beta-\delta,\max(\beta/2-\delta,1-\beta/2)\}
 \le \tfrac12+\tfrac\beta4-\tfrac\delta2.
\]
Every \(o(1)\) in this lemma is uniform and may be replaced
by an arbitrary fixed positive exponent.
The uncapped minimum without the entry \(1\) also gives a valid
counting exponent.  Its exact comparison with the affine expression
requires \(\beta-2\delta\le2\); the displayed capped comparison is
valid without this additional range restriction.
\end{lemma}

\begin{proof}
For a fixed \(q\), the equation \(n(n+h)=qv^2\) becomes
\((2n+h)^2-4qv^2=h^2\). Lemma~\ref{p1:arith:pell}
gives \(N^{o(1)}\) solutions, uniformly for the polynomial
heights in question. When \(q=1\), factor the difference of
squares instead. Summing over the \(O(1+N^\beta/e)\)
possible kernels gives the first exponent.

For the other estimate, \(v\asymp N^{1-\beta/2}\) and
\(ev^2\mid n(n+h)\). If
\(\rho_h(M)=\#\{a\bmod M:a(a+h)\equiv0\pmod M\}\), then
\[
 \rho_h(\ell^a)
 \le2\ell^{\min(v_\ell(h),\lfloor a/2\rfloor)},\qquad
 \rho_h(M)\le h\,2^{\omega(M)}.
\]
For \(a\le2v_\ell(h)\), a root is divisible by
\(\ell^{\lceil a/2\rceil}\); for larger \(a\), the two
separated roots give the displayed bound. The Chinese
remainder theorem proves the second inequality. Here
\(M=ev^2\le N^{2+o(1)}\), so summing residue classes gives
\[
 N^{o(1)}
 \sum_{v\asymp N^{1-\beta/2}}\left(\frac{N}{ev^2}+1\right)
 \le N^{o(1)}
       \{N^{\beta/2-\delta}+N^{1-\beta/2}\}.
\]
The trivial \(O(N)\) count supplies the entry \(1\) in the
minimum.  If \(\beta-\delta\ge1\), the maximum is
\(\beta/2-\delta\), and the capped exponent is at most
\[
 \min(1,\beta/2-\delta)
 \le \frac{1+\beta/2-\delta}{2}
 =\frac12+\frac\beta4-\frac\delta2.
\]
Otherwise the maximum is \(1-\beta/2\), and the minimum of this
number and \(\beta-\delta\) is at most their average, with the
same affine value.  This proves the exact comparison, including
admissible cases with \(\beta>2\).
\end{proof}

\paragraph{Necessity of the exponent cap.}
In Lemma~\ref{p13:short:quadratic}, the counting exponent is capped:
\[
 c_2(\beta,\delta)
   =\min\{1,\beta-\delta,\max(\beta/2-\delta,1-\beta/2)\}.
\]
Its affine upper bound is then valid for every real $\beta,\delta$.
The uncapped expression still gives a valid counting estimate, but
its exact affine comparison requires $\beta-2\delta\le2$ when
$\delta\ge0$. This restriction is necessary on actual positive-divisor
data: squarefree consecutive $N,N+1$ with $h=e=1$ give
$\beta=\log(N(N+1))/\log N>2$ and $\delta=0$.
The extra entry $1$ is supplied by the trivial $O(N)$ count.


\begin{lemma}[Kernels of a split product]
\label{p13:short:split}
Let
\[
 P_A(X)=\prod_{a\in A}(X+a),\qquad
 A\subset[-DB,DB]\cap\Z,\qquad
 2\le d=|A|\le c\log\log N,
\]
where \(d\) is even and \(D,c>0\) are fixed.
For a fixed squarefree \(e\), the number of \(n\asymp N\)
with \(q=\sf(P_A(n))\asymp N^\beta\) and \(e\mid q\) is at most
\begin{equation}
 N^{1/2+\beta/(2d)-\delta/d+c\log2+o(1)},
 \qquad \delta=\frac{\log e}{\log N}.
 \label{p13:short:split-count}
\end{equation}
For a fixed kernel \(q\), the fibre has at most
\(N^{2c\log2+o(1)}\) elements. All estimates are uniform in \(A\).
\end{lemma}

\begin{proof}
Pair the \(d\) linear factors into \(m=d/2\) quadratics
\(Q_j\), put \(q_j=\sf(Q_j(n))\), and write
\[
 D_A=\prod_{\substack{a,b\in A\\a<b}}|a-b|.
\]
Its logarithm is \(O_D(d^2\log B)=o(\log N)\).
Outside \(D_A\), the \(q_j\) are pairwise coprime, and
\begin{equation}
 \frac{\prod_jq_j}{q}\mid D_A^2.
 \label{p13:short:resultant}
\end{equation}
Indeed, a prime occurring in \(s\) of the \(q_j\) divides at
least \(s\) distinct linear factors. Its exponent in the
quotient is \(2\lfloor s/2\rfloor\), whereas its exponent
in \(D_A^2\) is at least \(2\binom{s}{2}\).

Remove the primes of \(D_A\) from \(e\), obtaining
\(e'=e/\gcd(e,\operatorname{rad}D_A)\).
There is a factorization \(e'=\prod_j e_j\), with
\(e_j\mid q_j\), such that the factors are pairwise coprime.
Set \(\beta_j=\log q_j/\log N\) and
\(\delta_j=\log e_j/\log N\). Then
\[
 \sum_j\beta_j=\beta+o(1),\qquad
 \sum_j\delta_j=\delta+o(1).
\]
For at least one \(j\),
\(\beta_j-2\delta_j\le(\beta-2\delta)/m+o(1)\).
Lemma~\ref{p13:short:quadratic} bounds its quadratic count
by \(N^{1/2+\beta/(2d)-\delta/d+o(1)}\).
We enumerate only that \(j\), the divisor \(e_j\mid e'\),
and its dyadic size. This costs at most
\(m\tau(e')O(\log N)\), not the number of all factorizations
of \(e'\). Since \(e'\le(C_DN)^d\), the product of its distinct
prime factors is at least \(\omega(e')!\). Stirling's estimate
therefore gives
\[
 \omega(e')\le(c+o(1))\log N,\qquad
 \tau(e')\le N^{c\log2+o(1)}.
\]
This proves \eqref{p13:short:split-count}.

For a fixed \(q\), choose two distinct offsets \(a,b\).
Both \(\sf(n+a)\) and \(\sf(n+b)\) divide
\(\operatorname{rad}(qD_A)\), so there are at most
\(4^{\omega(qD_A)}\le N^{2c\log2+o(1)}\) choices.
The effective coefficients \(u,v\) satisfy \(u,v\ll_D N\).
For each such choice,
\(us^2-vt^2=a-b\ne0\) has \(N^{o(1)}\) solutions by
Lemma~\ref{p1:arith:pell}, including its factorizable case.
Thus Pell is applied to coefficients of height \(O(N)\),
even though \(q\) itself may be as large as \(N^{O(d)}\).
\end{proof}

\subsubsection{Proof of Proposition~\ref{p13:short:moments}}
\label{p14:comp:short:moments}

\begin{proof}
The moment counts the tuples for which
\(\prod_iP_A(n_i)\) is a square. Write \(q_i=\sf(P_A(n_i))\).
There is a unique representation
\[
 q_i=\prod_{\substack{S\subseteq\{1,\ldots,v\}\\
                  |S|\ge2\text{ even},\ i\in S}}d_S
\]
with pairwise coprime squarefree \(d_S\). Fix dyadic size
classes for these numbers, and put
\(\alpha_S=\log d_S/\log N\) and
\(\beta_i=\sum_{S\ni i}\alpha_S\).
There are \(N^{o(1)}\) size profiles, since \(v\) is fixed.

Expose \(v-1\) of the integers in a chosen order. Before
exposing \(i\), the product of the \(d_S\) whose sets meet
an earlier index is a prescribed divisor \(e_i\mid q_i\).
If its exponent is \(\delta_i\), Lemma~\ref{p13:short:split}
gives the cost
\[
 N^{1/2+\beta_i/(2d)-\delta_i/d+c\log2+o(1)}.
\]
After choosing \(n_i\), assign the new prime factors of
\(q_i/e_i\) to the at most \(2^{v-1}\) possible sets \(S\).
Their number is at most
\(N^{c(v-1)\log2+o(1)}\); inconsistent assignments are
discarded. Every \(S\) has at least two members, so the last
kernel is now fixed and its fibre costs
\(N^{2c\log2+o(1)}\). The enumeration costs in total at most
\(N^{c[v(v-1)+2]\log2+o(1)}\).

For the remaining exponent, average over orders. A set \(S\)
of size \(a\) has coefficient
\[
 \frac1d\left(1-\frac a2+
               \frac12\ind_{\{\text{last index}\in S\}}\right)
\]
in \(\sum_{i\ne\mathrm{last}}(\beta_i/(2d)-\delta_i/d)\).
Its average is \(d^{-1}(1-a/2+a/(2v))\), which is
\(1/(dv)\) for \(a=2\) and nonpositive for \(a\ge4\).
As \(2\sum_{|S|=2}\alpha_S\le\sum_i\beta_i\le vd+o(1)\),
some order has total exponent at most \(v/2+o(1)\).
The order can depend on the size profile; its \(v!\) choices
do not affect the estimate. This proves
\eqref{p13:short:moment-bound}.
\end{proof}
